\section{Conclusion}\label{sec:conclusion}

This dissertation partially reproduced selected PAC-Bayes-with-Backprop cells on
MNIST under a reduced-compute protocol, extended them to Fashion-MNIST and
CIFAR-10, and re-implemented the certificate-computation pipeline with corrected
bookkeeping.

On reproduction, the corrected pipeline matches the reference's MNIST stochastic
test errors to within $\ReproStochDeltaAbsMax$ on all four target cells, with
the sign of the difference varying between them. The certificates are uniformly
looser; the Monte-Carlo sweep accounts for roughly half of the gap on the cell
examined in detail, and the rest coincides with protocol differences, chiefly the
absent per-cell grid search, that we identified but did not measure. The
reproduction is therefore complete on the predictive metrics and partial on the
certificate. Decomposing it revealed effects that test error alone does not: it
separates into the empirical Gibbs risk and a complexity term scaling with
$\mathrm{KL}/n_{\mathrm b}$, and a learnt prior tightens it in every paired
comparison we ran, though not by the same route each time. Measured
on the scale that enters the bound, $\mathrm{KL}/n_{\mathrm b}$, the complexity
term is the larger route only on MNIST; on Fashion-MNIST the two routes are
comparable, and on CIFAR-10 the learnt prior pays \emph{more} complexity and wins
entirely on empirical risk. That decomposition is what makes the certificate
sensitive in the ways the sensitivity analysis probes --- to the prior, the
bound-set size, the Monte-Carlo budget and the seed --- and it is where the audit
proves its worth, since correcting the bound-set sample count alone loosens the
small-data certificates substantially. Difficulty was
therefore tested in the one design that controls for it: with everything but the
label noise held constant, the noise raises the empirical-risk term and, through
the PAC-Bayes-kl inversion, the certificate, so the effect can be attributed to
the perturbation itself, with the proviso that the certified quantity is then the
Gibbs risk on the noisy-label distribution.

Together these results show that PAC-Bayes-with-Backprop delivered formally
non-vacuous certificates in all $23$ reported cells, none reaching $1$ --- a
statement about runs that trained, since the $15$-layer cell is the mean of its
three converged seeds. Non-vacuous and useful are different thresholds, though:
the loosest, CIFAR-10 under $f_{\text{bbb}}$, is
$\CertCifaronezeroCnnBbbLearnt$, certifying almost nothing a practitioner would
act on, and tight certificates arrived only on the easier benchmarks. In all
three paired $f_{\text{quad}}$ controls the data-dependent prior produced a large
tightening --- not the largest effect in the study, since $40\%$ label noise
moves the same MNIST certificate about twice as far.
The value of any such number depends on the implementation that produced it, and
several conventions in the reference required correction in our
re-implementation. The main contribution of the project is that corrected, tested
pipeline: reliable certificates require correct sample counts, confidence
accounting, Monte-Carlo aggregation and seed handling, together with enough
provenance to re-derive each reported value.

The clearest direction for further work is to close the gap
between the certified Gibbs classifier and the deterministic predictor actually
deployed. The C-bound of \textcite{lacasse2006,germain2015} relates the
majority-vote risk to the mean and variance of the Gibbs risk, and adapting it to
the objectives used here would give a formal guarantee for a deterministic
predictor. Which one matters: the C-bound governs the majority vote induced by
$Q$, the limit of the ensemble predictor, not the posterior-mean network whose
error we report, so a guarantee for the latter would remain open. On
the empirical side, the full-budget re-draw reported here removes one identified
source of the discrepancy but not the absent grid search, which is what a
reference-matched hyper-parameter search would settle.
The CIFAR-10 results invite two extensions: the deeper 13- and 15-layer
architectures that the reference shows give tighter bounds, and the
differentially-private prior route, which would let the prior use all of the data
and is the more promising way to attack the high empirical risk that, as the
error analysis shows, is what keeps CIFAR-10 certificates loose. Finally, the
controlled-difficulty design could be broadened from label noise to input
corruption and class overlap, to test whether the empirical-risk mechanism
identified here is the dominant one across different notions of difficulty.
